\documentclass[11pt]{article}
\usepackage[utf8]{inputenc}

%----- Configuración del estilo del documento------%
\usepackage{epsfig,graphicx}
\usepackage[left=2cm,right=2cm,top=1.8cm,bottom=2.3cm]{geometry}
\usepackage{fancyhdr}
\usepackage{lastpage}
\usepackage{url}
\usepackage{longtable}
\pagestyle{fancy}
\fancyhf{}
\fancyhead[R]{\ifnum\value{page}>1 \textbf{C\'alculo 3}\fi}
\rfoot{\textit{Página \thepage}}

%Encabezado izquierdo: logo
\lfoot{\includegraphics[height=1.0cm]{Logo_FC.png}}



%------ Paquetes matemáticos básicos --------%
\usepackage{amsmath}
\usepackage{amssymb}
\usepackage{amsthm}
\usepackage[spanish]{babel}
\usepackage{graphicx}
\usepackage{hyperref}

\usepackage{tabularx}
\usepackage{xcolor}
\usepackage[table]{xcolor}
\usepackage{colortbl}
\usepackage{array, multirow, multicol, tabularx}
\usepackage{tcolorbox}
\newtheorem{theorem}{Theorem}[section]
\newtheorem{corollary}{Corollary}[theorem]
\newtheorem{lemma}[theorem]{Lemma}

%------si-------%
\definecolor{B}{HTML}{FFFFFF}
\definecolor{G}{HTML}{5e5e5e}
\definecolor{R2}{HTML}{d53d40}
\definecolor{A2}{HTML}{034190}
\definecolor{V2}{HTML}{7faa50}

\newcommand{\PrdI}[2]{\ensuremath{\langle #1, #2 \rangle}}
\newcommand{\R}{\mathbb{R}}
\newcommand{\C}{\mathbb{C}}
\newcommand{\N}{\mathbb{N}}
\newcommand{\Z}{\mathbb{Z}}
\newcommand{\Q}{\mathbb{Q}}
\renewcommand{\theenumi}{\Roman{enumi}}
\renewcommand{\labelenumi}{{\theenumi}.}

\begin{document}
	\makeatletter
        \renewenvironment{proof}[1][\proofname]{\par
            \pushQED{\qed}%
            \normalfont \topsep6\p@\@plus6\p@\relax
            \trivlist
            \item\relax
            {\itshape
            #1\@addpunct{.}}\par\vspace{\baselineskip}\ignorespaces
            }{%
            \popQED\endtrivlist\@endpefalse
            }
        \newenvironment{Solucion}[1][Soluci\'on]{\par
            \pushQED{\qed}%
            \normalfont \topsep6\p@\@plus6\p@\relax
            \trivlist
            \item\relax
            {\itshape
            #1\@addpunct{.}}\par\vspace{\baselineskip}\ignorespaces
            }{%
            \popQED\endtrivlist\@endpefalse
            }
            
    \makeatother

%------ Encabezado -------- %

\begin{center}
    \begin{minipage}{3cm}
    	\begin{center}
    		\includegraphics[height=3.4cm]{logo_unam.png}
    	\end{center}
    \end{minipage}\hfill
    \begin{minipage}{10cm}
    	\begin{center}
    	\textbf{\large Universidad Nacional Autónoma de México}\\[0.1cm]
        \textbf{Facultad de Ciencias}\\[0.1cm]
        \textbf{Ejercicios }\\[0.1cm]
         C\'alculo 3 \\[0.1cm]
         El\'ias L\'opez Rivera\\[0.1cm]
        \texttt{elias.lopezr,@ciencias.unam.mx }\\[0.1cm]
        Fecha:\,\,28/12/2025
    	\end{center}
    \end{minipage}\hfill
    \begin{minipage}{3cm}
    	\begin{center}
    		\includegraphics[height=3.4cm]{Logo_FC.png}
    	\end{center}
    \end{minipage}
\end{center}
\rule{17cm}{0.1mm}
\begin{tcolorbox}[title = \textcolor{black}{\textcolor{white}{Problema 1}}]
\textit{Sea $\{F_{n}\}_{n\in \N}$ es una sucesi\'on anidada $(F_{n+1}\subseteq F_n\,\,\forall n\in \N)$ de conjuntos compactos no vacios entonces:\,\\
\begin{equation*}
    \bigcap_{n=1}^{\infty}F_n\neq \emptyset
\end{equation*}}
\end{tcolorbox}
\begin{proof}
    Existe una prueba para este hecho sin necesidad de usar la definici\'on de compacidad, es bastante ilustrativa de algunos trucos clasicos en an\'alisis, se le recomienda ampliamente al lector consultarla o intertarla por su cuenta,
    en el presente texto no la usaremos.Procedemos por reducci\'on al absurdo, supongamos que:\,\\
    \begin{equation*}
        \bigcap_{n=1}^{\infty}F_n= \emptyset\implies \bigcup_{n=1}^{\infty}F_n^c=\R^n
    \end{equation*}\,\\
    luego $F_1$ es compacto y $F_1\subseteq \R^n=\bigcup_{n=1}^{\infty}F_n^c$, por lo que existe una subcubierta abierta finita, digamos $\displaystyle\bigcup_{n\in \N_k}F^c_{n}$ tal que:\,\\
    \begin{equation*}
        F_{k}\subseteq F_1\subseteq \bigcup_{n\in \N_k}F_n^c
    \end{equation*}\,\\
    Luego notemos que si $k\geq n$ entones $F_{k}\subset F_n$, que es equivalente a que $F_{k}\cap F^c_n=\emptyset$, por tanto:\,\\
    \begin{equation*}
       F_k=F_{k}\cap \bigcup_{n\in \N_k}F_n^c=\bigcup_{n\in \N_k}(F_{k}\cap F_n^c)=\emptyset
    \end{equation*}\,\\
    lo cual es una contradicci\'on pues $F_k$ es no vacio por hip\'otesis, con lo que el resultado queda demostrado.
\end{proof}
\begin{tcolorbox}[title = \textcolor{black}{\textcolor{white}{Problema 1}}]
\textit{Sean $g:V\subset \R^n\to \R^n$ derivable en $y_0\in V$ tal que  $D_g(y_0)$ es invertible, y sea $f:U\subset \R^n\to \R^n$ tal que 
$f(U)\subset V$ y $x_0\in U$ tal que $f(x_0)=y_0$.Pruebe que si $g\circ f$ es derivable en $x_0$, y existen $0<M$ y $0<r$ tales que:\,\\
\begin{equation*}
    \frac{\|f(x)-f(y)\|}{\|x-y\|}\leq M
\end{equation*}\,\\
para toda $x\in B_{r}\setminus \{x_0\}$, entonces:
\begin{enumerate}
   \item $f$ es derivable en $x_0$ 
   \item Encuentre una expresi\'on para $D_f(x_0)$ en t\'erminos de $D_{g\circ f}(x_0)$ y $D_g(y_0)$
\end{enumerate}}
\end{tcolorbox}
\begin{proof}
\end{proof}
\end{document}

